% from ega0-18-fr.tex
\oldpage[0\textsubscript{IV-1}]{62}

\begin{env}[18.3.6.4]
\label{0.18.3.6.4-fr}
Enfin, l'application $(u,1_B,1_L)_*$ est l'application
\[
  u^*:\operatorname{Exan}_A(B,L)\longrightarrow
  \operatorname{Exan}_{A'}(B,L)
\]
qui s'obtient en faisant correspondre à une $A$-extension $E$ de $B$ par $L$
l'anneau $E$ considéré comme $A'$-\emph{anneau} au moyen de $u$ (18.1.1), qui
est évidemment une $A'$-extension de $B$ par $L$, $B$ étant aussi considéré
comme $A'$-anneau au moyen de $u$~; il est clair qu'une $A$-équivalence est
aussi une $A'$-équivalence, d'où l'application (18.3.6.4), qui, pour tout
homomorphisme $w:L\to L'$ de $B$-bimodules, rend encore commutatif le diagramme
\begin{equation}
\label{0.18.3.6.5-fr}
\tag{18.3.6.5}
  \xymatrix@C=3pc@R=2.4pc{
    \operatorname{Exan}_A(B,L)\ar[r]^{u^*}\ar[d]_{w_*}
      & \operatorname{Exan}_{A'}(B,L)\ar[d]^{w_*}\\
    \operatorname{Exan}_A(B,L')\ar[r]_{u^*}
      & \operatorname{Exan}_{A'}(B,L')
  }
\end{equation}
d'où l'on conclut comme ci-dessus que $u^*$ est un homomorphisme de groupes.

Cela étant, on pose
$(u,v,w)_*=(1_{A'},1_{B'},w)_*\circ(1_{A'},v,1_L)_*
\circ(u,1_B,1_L)_*$ et l'on vérifie facilement, en raison de la commutativité
des diagrammes (18.3.6.3) et (18.3.6.5), que l'on a
\[
  (u\circ u',v\circ v',w'\circ w)_*
  =(u',v',w')_*\circ(u,v,w)_*,
\]
ce qui achève de prouver que (18.3.4.1) est un foncteur.

L'existence de l'homomorphisme de groupes (18.3.6.1) montre en particulier
que si $E$ est une $A$-extension \emph{triviale} de $B$ par $L$, l'extension
$E\oplus_LL'$ de $B$ par $L'$ définie dans (18.2.8) est aussi triviale, ce que
l'on vérifie d'ailleurs sans peine de façon directe.

D'autre part, pour tout élément $z$ du \emph{centre} $Z$ de $B$, l'homothétie
$h_z:y\mapsto yz=zy$ est un endomorphisme du $B$-bimodule $L$, donc $(h_z)_*$
est un endomorphisme du groupe commutatif $\operatorname{Exan}_A(B,L)$, et par
fonctorialité, ces endomorphismes définissent sur
$\operatorname{Exan}_A(B,L)$ une structure canonique de $Z$-\emph{module}.
\end{env}

\begin{env}[18.3.7]
\label{0.18.3.7-fr}
Soient $A$, $A'$ deux anneaux, $u:A'\to A$ un homomorphisme, $B$ un
$A$-anneau et $L$ un $B$-bimodule. Le \emph{noyau} de l'homomorphisme de
groupes
\[
  u^*:\operatorname{Exan}_A(B,L)\longrightarrow
  \operatorname{Exan}_{A'}(B,L)
\]
est formé par définition des classes de $A$-extensions de $B$ par $L$ qui
sont $A'$-\emph{triviales} quand on les considère comme $A'$-extensions au
moyen de $u$. On désigne ce noyau par la notation
$\operatorname{Exan}_{A/A'}(B,L)$ quand cela ne prête pas à confusion.

Si $\Lambda$ est un anneau, et si on a un diagramme commutatif de
$\Lambda$-\emph{homomorphismes} de $\Lambda$-anneaux
\begin{equation}
\label{0.18.3.7.1-fr}
\tag{18.3.7.1}
  \xymatrix@C=2pc@R=1.5pc{
    B'\ar[r] & B\\
    A'\ar[u]\ar[r] & A\ar[u]
  }
\end{equation}

